\documentclass[10pt,a4paper]{amsart}
\usepackage[margin=19mm]{geometry}
\usepackage{amsmath,amssymb,mathtools}
\usepackage[scr=rsfs,cal=euler]{mathalfa}
\usepackage{xfrac}
\usepackage[unicode,hidelinks,bookmarks=false]{hyperref}
\newcommand{\ie}{i.e.\ }
\newcommand{\cf}{cf.\ }
\newcommand{\resp}{respectively\ }
\newcommand{\Subsection}{\subsection}
\providecommand{\nobreakdash}{\nobreak\hbox{-}}
\setlength{\emergencystretch}{2em}
\multlinegap0pt
\usepackage{mathtools}
\usepackage{amssymb}
\usepackage{siunitx}
\usepackage[shortlabels]{enumitem}
\usepackage{bm}
\newtheorem{theorem}{Theorem}
\newtheorem{lemma}{Lemma}
\newtheorem{defn}{Definition}
\title{\vspace{-0.35em}A logical treatment of noise:\\ Solving The Hardest Logic Puzzle Ever and its generalizations}
\author{Daniel Vallstrom}
\date{Source: arXiv:2201.09801v12 (CC BY 4.0)}
\begin{document}
\maketitle

\noindent\emph{Source and licence.} \vspace{-0.35em}A logical treatment of noise:\\ Solving The Hardest Logic Puzzle Ever and its generalizations. Version arXiv:2201.09801v12 (CC BY 4.0), distributed by arXiv under the Creative Commons Attribution 4.0 International licence, http://creativecommons.org/licenses/by/4.0/, as recorded in the arXiv licence metadata for this exact version and shown on the abstract page https://arxiv.org/abs/2201.09801v12. Full licence notice: \texttt{licenses/arxiv-2201.09801-CC-BY-4.0.txt}.\par\smallskip

\noindent\textit{Editorial scope.} Counted unit: the article's theorem theorem:nGodsSolvability (source0.tex lines 672--676). The article's complete original proof of it is delivered with it (source0.tex lines 728--925, one \texttt{proof} environment of 198 lines, which the article places away from the statement). No mathematics is written, repaired or re-derived: the statement and the proof are the article's own lines, in the article's own notation, and no retained line is substituted or re-flowed.  Retained uncounted as the source-local context the proof needs: lines 714--718.  Nothing was omitted from any retained span, so the package carries no omission record.  No retained line contains a citation, so the wrapper carries no bibliography and no bibliography entry.  No retained line contains a figure environment, an image-inclusion command or drawing code, so no image is delivered and the package declares no asset dependency.  Editorial formatting only, in addition: an AMS \texttt{amsart} preamble that restates the article's own declarations and macros used by the retained lines, namely the environments \texttt{theorem}, \texttt{lemma}, \texttt{defn} on separate counters, together with the standard packages those lines need.  The retained environments print in this document as Theorem 1, Lemma 1, Definition 1 in order of appearance, whereas the article prints the same items with the numbering of its own full document.  The complete native source of the article as distributed in the arXiv e-print for this exact version is bundled unchanged as \texttt{sources/121150/source0.tex}.
\begin{lemma}
\label{lemma:RsLTNonRsSolvable}
If an $n$ gods puzzle has strictly more non-random gods than random gods, 
then it is solvable.  
\end{lemma}

\begin{theorem}
\label{theorem:nGodsSolvability}
An $n$ gods puzzle is solvable if and only if the number of random gods is strictly less than
the number of non-random gods.
\end{theorem}

\begin{proof}[Proof of theorem \ref{theorem:nGodsSolvability}]
The $\Leftarrow$ case is covered by lemma \ref{lemma:RsLTNonRsSolvable}.

For the $\Rightarrow$ case,
to see why puzzles with at least as many random gods 
as non-random gods aren't solvable, consider the easiest to solve of such problems:
Assume that $n$ is even, that the random gods equal the non-random gods,
and that there are only truthful non-random gods.
Let 
\begin{defn}
\label{def:pnUns}
$p_n \coloneqq $
 \begin{equation}
 \label{eq:nonSolvable}
  \bigvee\!\!
   \left\lbrace
    \begin{matrix}
     \land\{\gamma_1\!=\!\mathcal{T},\gamma_2\!=\!\mathcal{R},\gamma_3\!=\!\mathcal{T},
     \gamma_4\!=\!\mathcal{R},\ldots,\gamma_n\!=\!\mathcal{R}\},\\
     \land\{\gamma_1\!=\!\mathcal{R},\gamma_2\!=\!\mathcal{T},\gamma_3\!=\!\mathcal{R},
     \gamma_4\!=\!\mathcal{T},\ldots,\gamma_n\!=\!\mathcal{T}\}\;\\
    \end{matrix}\!
   \right\rbrace\!\!
 \end{equation}
 % This eq. is enough; we don't have to add that randoms are equal or more than
 % non-randoms. E.g.\ n=3 is fine, as long as both conjunctions can hold.
\end{defn}
Let $p_{\mathcal{\scriptscriptstyle T}\mathcal{\scriptscriptstyle R}}$ and
$p_{\mathcal{\scriptscriptstyle R}\mathcal{\scriptscriptstyle T}}$ be the
conjunctions of $p_n$:
\begin{defn}
\label{pTR}
$p_{\mathcal{\scriptscriptstyle T}\mathcal{\scriptscriptstyle R}}\coloneqq $
 \begin{equation}
   \land\{\gamma_1\!=\!\mathcal{T},\gamma_2\!=\!\mathcal{R},\gamma_3\!=\!\mathcal{T},
   \gamma_4\!=\!\mathcal{R},\ldots,\gamma_n\!=\!\mathcal{R}\}
 \end{equation}
\end{defn}
\begin{defn}
\label{pRT}
$p_{\mathcal{\scriptscriptstyle R}\mathcal{\scriptscriptstyle T}}\coloneqq $
 \begin{equation}
  \hspace{3.0em}\land\{\gamma_1\!=\!\mathcal{R},\gamma_2\!=\!\mathcal{T},
  \gamma_3\!=\!\mathcal{R},
  \gamma_4\!=\!\mathcal{T},\ldots,\gamma_n\!=\!\mathcal{T}\}
 \end{equation}
\end{defn}
If the gods are as described by 
$p_{\mathcal{\scriptscriptstyle T}\mathcal{\scriptscriptstyle R}}$ or
$p_{\mathcal{\scriptscriptstyle R}\mathcal{\scriptscriptstyle T}}$, and only
$p_n$ is known, then that puzzle instance is unsolvable, if the random gods are 
unhelpful (we'll show this next).
So let's assume that the gods are as described by $p_n$.

To get to an unsolvable position, we'll have the random gods ``happen''
to force the puzzle there. To that end,
%given already known conclusions $Q$,
if both $q$ and $\neg q$ are consistent with $p_n$,
(i.e.\ not $p_n\!\rightarrow\!\neg q$ and not $p_n\!\rightarrow\!q$),
assume that the random gods always happen to give the incorrect answer
to a question about $q$.
If only one of $q$ and $\neg q$ is consistent with $p_n$,
assume that the random gods always happen to give the correct answer 
to a question about $q$.

%  ---
% on questions that doesn't follow yet from what is known that is. 
% For already known answers they should happen to answer non-contradictory. 
% For example, they shouldn't say that $\bot$ is true.

Call a question `trivial' if it or its negation follow from already
asked questions.

Given this setup,
and if non-trivial questions are asked,
then it will be known that $p_n$ holds.
But once there,
no more conclusion can be drawn from the answer to a question $q$.
And it is not possible to determine which of 
$p_{\mathcal{\scriptscriptstyle T}\mathcal{\scriptscriptstyle R}}$ and
$p_{\mathcal{\scriptscriptstyle R}\mathcal{\scriptscriptstyle T}}$
that holds.

More specifically, without loss of generality we will assume that $\gamma_1$ is asked about $q$. 

% To see that $p_n$ will be concluded if non-trivial questions are asked,
% note that what we are able to conclude are
% $q \vee \gamma_1\!=\!\mathcal{R}$ from $t(q,\gamma_1)$, and
% $\neg q \vee \gamma_1\!=\!\mathcal{R}$ from $\neg t(q,\gamma_1)$. 
% Since $\gamma_1\!=\!\mathcal{R}$ is consistent with $p_n$, any disjunction with it is too.
% (Alternatively, 
% if $\gamma_1\!=\!\mathcal{T}$, then an answer that $q$ (or $\neg q$) holds
% implies that $q$ (or $\neg q$) must be consistent with $p_n$ since $p_n$ is in fact true,
% and $\gamma_1$ is non-random. 
% %Besides, $\gamma_1\!=\!\mathcal{R}$ too is consistent with $p_n$.
% And if $\gamma_1\!=\!\mathcal{R}$, then again $q$ (or $\neg q$) is consistent with $p_n$  
% because of the assumptions of how random gods answer.)
% %1) q,-q both consi
% %2) only one consi

%Eventually $p_n$ will be concluded from any ordinary non-trivial questions. However, 
%the easiest way to 
An easy way to see that we 
%can
are able to
conclude $p_n$ is to note that we can ask
%If nothing else, eventually
$\frac{n}{2}+1$ gods
% are asked 
about $p_n$ explicitly.
Since all gods answer that $p_n$ holds, $p_n$ can be concluded since at least
$1$ of $\frac{n}{2}+1$ gods must be non-random.

% we are able to conclude only
% $q \vee \gamma_1\!=\!\mathcal{R}$ (or $\neg q \vee \gamma_1\!=\!\mathcal{R}$).
% And the false answer that $q$ (or $\neg q$) holds
% undoes the disjunct $\gamma_1=\mathcal{R}$. 

Note that
if $p_n$ 
%holds
is known
and $q$ is non-trivial, then
$q\leftrightarrow p_{\mathcal{\scriptscriptstyle T}\mathcal{\scriptscriptstyle R}}$ or
$q\leftrightarrow p_{\mathcal{\scriptscriptstyle R}\mathcal{\scriptscriptstyle T}}$.
Because suppose 
that $q$ is non-trivial. Then it has a model $\mathcal{M}$ (i.e.\ an assignment of the gods)
where $q$, and $p_n$, are true. Suppose, without loss of generality, that it's 
the $p_{\mathcal{\scriptscriptstyle T}\mathcal{\scriptscriptstyle R}}$
disjunct that's true in $\mathcal{M}$.
Then, since $p_{\mathcal{\scriptscriptstyle T}\mathcal{\scriptscriptstyle R}}$
completely determines a model, $q$ holds whenever 
$p_{\mathcal{\scriptscriptstyle T}\mathcal{\scriptscriptstyle R}}$ does, i.e.\
$p_{\mathcal{\scriptscriptstyle T}\mathcal{\scriptscriptstyle R}}\rightarrow q$.
Suppose $\neg p_{\mathcal{\scriptscriptstyle T}\mathcal{\scriptscriptstyle R}}$.
Then $p_{\mathcal{\scriptscriptstyle R}\mathcal{\scriptscriptstyle T}}$.
Since $p_{\mathcal{\scriptscriptstyle R}\mathcal{\scriptscriptstyle T}}$
determines its models completely, if $q$ where to hold, then $q$ would
hold in all models to $p_n$, contradicting that $q$ is non-trivial.
Hence $\neg q$ holds, i.e.\ 
$\neg p_{\mathcal{\scriptscriptstyle T}\mathcal{\scriptscriptstyle R}}\rightarrow \neg q$.

% otherwise. Then there are models 
% (i.e.\ assignments of the gods)
% such that 
% $q \wedge p_{\mathcal{\scriptscriptstyle T}\mathcal{\scriptscriptstyle R}}$ is true in one and
% $q \wedge p_{\mathcal{\scriptscriptstyle R}\mathcal{\scriptscriptstyle T}}$ is true in an other.
% Since $p_n$ holds, one of its disjuncts must hold.
% Since 
% $p_{\mathcal{\scriptscriptstyle T}\mathcal{\scriptscriptstyle R}}$ and
% $p_{\mathcal{\scriptscriptstyle R}\mathcal{\scriptscriptstyle T}}$
% completely determine a model, 
% $p_{\mathcal{\scriptscriptstyle T}\mathcal{\scriptscriptstyle R}}$ and
% $p_{\mathcal{\scriptscriptstyle R}\mathcal{\scriptscriptstyle T}}$
% desribe all of the models of $p_n$. Hence $q$ is true in all models of $p_n$,
% contradicting the assumption that $q$ is non-trivial.

Assume that $p_n$ is known.
To see that once $p_n$ is known, no more conclusions can be made, 
suppose $\gamma_1\!=\!\mathcal{T}$, 
and $q$ is non-trivial.
Then an answer that $q$ (or $\neg q$) holds 
(with implications that 
$p_{\mathcal{\scriptscriptstyle T}\mathcal{\scriptscriptstyle R}}$ holds)
is undone because we can only conclude that
$q \vee \gamma_1\!=\!\mathcal{R}$ (or $\neg q \vee \gamma_1\!=\!\mathcal{R}$),
which is equivalent to $p_n$.

Similarly, if instead $\gamma_1=\mathcal{R}$, then we are able to conclude only
$q \vee \gamma_1\!=\!\mathcal{R}$ (or $\neg q \vee \gamma_1\!=\!\mathcal{R}$).
If $q$ is non-trivial, 
then both $q$ and $\neg q$ are consistent with $p_n$.
Hence, the false answer that $q$ (or $\neg q$) holds
undoes the disjunct $\gamma_1=\mathcal{R}$. 

If $q$ is trivial, the concluded disjunction, 
$q \vee \gamma_1\!=\!\mathcal{R}$ (or $\neg q \vee \gamma_1\!=\!\mathcal{R}$),
will already be known, and nothing new can be concluded again.

% More formally, if $q$ is non-trivial, then both $q$ and $\neg q$ are consistent. But in their
% models, due to eq.\,\eqref{eq:nonSolvable}, everything is determined. Hence, no more conclusions
% can be made from any answer.
% This works for when eq. 8 isn't known yet to solver. g1=T case works because g1=R is always
% added as disjunct.
% g1=R works because that's how R answers, we'll assume. R lies only w.r.t. eq 8, not what's 
% known to solver. if q isn't consistent with eq 8, don't answer/lie that it holds.
% So situation will move towards eq. 8, and then stay there.

To see that even harder puzzles to solve are unsolvable too,
the restriction that the non-random gods are only truthful is immaterial.

If more random gods are added, the proof still works with minor alterations, e.g.\
to eq.\,\eqref{eq:nonSolvable}. 
% If there are more random than non-random gods, the problem can be reduced to
% a problem with equally many random as non-random gods.
% (Alternatively, it ought to be possible to reduce the
% problem to one where the random and non-random gods are equally many.)
\end{proof}

\end{document}
